\documentclass[10pt,a4paper]{amsart}
\usepackage[margin=19mm]{geometry}
\usepackage{amsmath,amssymb,mathtools}
\usepackage[scr=rsfs,cal=euler]{mathalfa}
\usepackage{xfrac}
\usepackage[unicode,hidelinks,bookmarks=false]{hyperref}
\newcommand{\ie}{i.e.\ }
\newcommand{\cf}{cf.\ }
\newcommand{\resp}{respectively\ }
\newcommand{\Subsection}{\subsection}
\providecommand{\nobreakdash}{\nobreak\hbox{-}}
\setlength{\emergencystretch}{2em}
\multlinegap0pt
\PassOptionsToPackage{normalem}{ulem}
\usepackage{ulem}
\usepackage[authoryear]{natbib}
\bibpunct[, ]{[}{]}{,}{a}{,}{}
\makeatletter
\numberwithin{equation}{section}
\numberwithin{figure}{section}
\theoremstyle{plain}
\newtheorem{thm}{\protect\theoremname}
\theoremstyle{plain}
\newtheorem{fact}[thm]{\protect\factname}
\newenvironment{lyxlist}[1]
	{\begin{list}{}
		{\settowidth{\labelwidth}{#1}
		 \setlength{\leftmargin}{\labelwidth}
		 \addtolength{\leftmargin}{\labelsep}
		 \renewcommand{\makelabel}[1]{##1\hfil}}}
	{\end{list}}
\theoremstyle{plain}
\newtheorem{cor}[thm]{\protect\corollaryname}
\theoremstyle{plain}
\newtheorem{prop}[thm]{\protect\propositionname}
\theoremstyle{remark}
\newtheorem*{claim*}{\protect\claimname}
\theoremstyle{plain}
\newtheorem{lem}[thm]{\protect\lemmaname}
\theoremstyle{definition}
\newtheorem{defn}[thm]{\protect\definitionname}
\makeatother
\providecommand{\claimname}{Claim}
\providecommand{\corollaryname}{Corollary}
\providecommand{\definitionname}{Definition}
\providecommand{\factname}{Fact}
\providecommand{\lemmaname}{Lemma}
\providecommand{\propositionname}{Proposition}
\providecommand{\theoremname}{Theorem}
\title[Internal Categoricity and the Generic Multiverse]{Internal Categoricity and the Generic Multiverse: the consistency of a non-standard model of $MV$ all of whose worlds are uncountable and transitive (Proposition 10)}
\author{Toby Meadows}
\date{Source: arXiv:2508.21202v1 (CC BY 4.0)}
\begin{document}
\begin{abstract}
John Steel's theory, $MV$, of the generic multiverse provides a foundation
for mathematics that aims to neutralize the effects of incompleteness
brought on by forcing arguments. Jouko V\"{a}\"{a}n\"{a}nen's development of internal
categoricity arguments provides opportunities to argue that the subject
matter of some theory is, in some sense, determined. This paper investigates
whether $MV$ is internally categorical.
\end{abstract}
\maketitle
\noindent\emph{Source and licence.} Internal Categoricity and the Generic Multiverse. Version arXiv:2508.21202v1 (CC BY 4.0), distributed under the Creative Commons Attribution 4.0 International licence, http://creativecommons.org/licenses/by/4.0/, as recorded in the arXiv OAI licence metadata for this exact version and shown on the abstract page https://arxiv.org/abs/2508.21202v1. Full rights notice: licenses/arxiv-2508.21202-CC-BY-4.0.txt.\par\smallskip
\noindent\textit{Editorial scope.} Counted unit: the unlabelled proposition asserting that if $L_{\omega_{1}}\models ZFC$ then it is consistent that there is a model $\mathcal{W}$ of $MV$ where every world $M$ in $\mathcal{W}$ is uncountable and transitive but $\mathcal{W}$ is not standard (source0.tex lines 780--784), printed as Proposition 10 in the published version, delivered with its complete original author proof at source0.tex lines 786--902 (117 lines, the longest proof block of the paper) as the single primary proof body. The proof is a standard \textbackslash{}begin\{proof\}/\textbackslash{}end\{proof\} block opening on the line after the statement, so it is adjacent and not detached; the paper's one detached proof block (source0.tex lines 608--619, which proves thm:MV|-ZFC\_count and fact:Steel comp) is retained inside the uncounted context and is NOT the counted proof. The paper's own section 2, \emph{What is $MV$?} (source0.tex lines 345--779), is retained verbatim and uncounted as the complete source-local core: the notation, the statement of the axioms of $MV$ with the discussion of their motivation, the subsection on the $MV$ theory with Fact 3, Theorem 4 and the proof of $MV\vdash ZFC^{-}_{count}$, the subsection on the models of $MV$ with Fact 5, Corollary 6, the proof of Corollary 6, Proposition 7, Corollary 8 and Proposition 9, and the paragraph that announces the counted proposition. Every reference of the retained spans resolves inside the unit, so no reference needed adaptation and the delivery evidence carries no replacement: the \textbackslash{}ref uses of the retained spans point only to labels declared inside them (cor:Steel, fact:Steel comp, thm:MV|-ZFC\_count, prop:Hjorth). The paper is itself an amsart article (source0.tex line 4: \textbackslash{}documentclass[oneside,english]\{amsart\}), so no publisher class is replaced, shipped or redistributed; the wrapper is the lane's pinned amsart editorial wrapper, and the paper's own settings are copied verbatim from the source: its theorem declarations as written by LyX (source0.tex lines 26--50, that is \textbackslash{}numberwithin for equation and figure, the plain, remark and definition theorem styles, \textbackslash{}newtheorem for thm, fact, cor, prop, lem and defn sharing the single thm counter, and the starred claim*), its lyxlist environment (source0.tex lines 34--40), its environment names (source0.tex lines 67--73), its \textbackslash{}usepackage[authoryear]\{natbib\} (source0.tex line 17) and its bibliography helper declarations \textbackslash{}natexlab, \textbackslash{}url and \textbackslash{}doi (source0.tex lines 2057--2061). One typographic declaration of the wrapper is disclosed: the e-print source carries no \textbackslash{}bibpunct, yet the published PDF prints citations in square brackets, so the wrapper adds \textbackslash{}bibpunct[,]\{[\}\{]\}\{,\}\{a\}\{,\}\{\} in order that the delivered citations reproduce the punctuation of the version of record ([Author, Year] for \textbackslash{}citep and Author [Year] for \textbackslash{}citet); it changes no number, symbol, statement or proof step. The paper's other preamble loads are not reproduced because nothing in the retained spans uses them: ae, aecompl, fontenc, inputenc, babel, geometry, setspace and ulem are font, encoding, page-geometry and spacing matter supplied by the wrapper, while qtree, bussproofs, pdflscape, tikz and tikz-cd are used only outside the retained spans (the paper has 38 tikz-related lines and one qtree line, none of them between source0.tex lines 345 and 902, and the retained spans contain no tikzpicture, no prooftree and no landscape environment). Every printed number of the delivered unit is produced by the paper's own counters: the wrapper is initialized editorially with \textbackslash{}setcounter\{section\}\{1\} and \textbackslash{}setcounter\{thm\}\{2\} immediately before the retained span, because the only numbered theorem-like environments preceding it are Theorem 1 and Theorem 2 of section 1 and the only section preceding it is section 1, so that the paper's own section heading prints as 2, and Fact 3, Theorem 4, Fact 5, Corollary 6, Proposition 7, Corollary 8, Proposition 9 and the counted Proposition 10 print exactly as in the published version (the paper numbers theorem-like environments consecutively through the article, without a per-section reset), with subsections 2.1, 2.2 and 2.3 and the paper's own \textbackslash{}numberwithin\{equation\}\{section\}; the delivered spans contain only unnumbered displays (align* and \textbackslash{}[ \textbackslash{}]), and I verified in the published PDF that Proposition 9 is printed for prop:Hjorth and Proposition 10 for the counted statement. Excluded and uncounted so that no other material inflates difficulty: the LyX front matter and the introduction (source0.tex lines 1--344 and 93--344, of which only the paper's own abstract is reproduced verbatim as front matter), the rest of the paper from source0.tex line 903 to the end (Proposition 11, all of section 3, sections 4 and 5 with the internal categoricity and solidity theorems, the discussion and the conclusion), and the paper's bibliography listing. The retained spans contain no figure, no table and no artwork: the paper has no includegraphics, no figure environment and no external graphic file, and its two tikzpictures lie outside the retained spans, so no delivered span has an includegraphics. The paper's source contains no non-ASCII character at all (verified count zero), so no missing glyph can arise from the retained text. Citations: the paper cites with the natbib author-year commands \textbackslash{}citep and \textbackslash{}citet, so the delivered bibliography is the paper's own in-source \textbackslash{}begin\{thebibliography\} block (source0.tex lines 2056--2321, 45 entries) transcribed verbatim with its own header declarations, reduced to exactly the 16 keys cited inside the retained spans, namely WoodinGM, SteelGP, MeadMadPGMV, Meadows2Args, BAGARIA\_TERNULLO\_2023, BLUE\_2024, JechST, KunenST2, foreman2009handbook, lindstrm2003aspects, EnayatLelykFOCat, ReitzTGA, Cohen1963-COHTIO-5, Solovay1970AMO, HjorthU2 and KunenST, each with the paper's own printed author-year label, so every citation of the delivered unit is the published one; the remaining 29 uncited entries of the paper are omitted. No mathematics is written, repaired or re-derived: the statement, the proof and the retained context are the paper's own text line for line, including the author's own footnotes, his own claim* inside the proof, his own metamathematical asides and his own printed slips.
\setcounter{section}{1}
\setcounter{thm}{2}
\section{What is $MV$?\label{sec:What-is-?}}

The incompleteness of $ZFC$ as evidenced by forcing arguments is
sometimes thought of as particularly deep. By contrast, shallow forms
of incompleteness can be resolved by adopting a stronger theory that
we think is true. For example, although $ZFC$ cannot prove that it
is consistent, if we adopt an axiom saying that there are inaccessible
cardinals then the consistency of $ZFC$ follows almost immediately.
But over at the deep end, the continuum hypothesis can be both forced
to hold and forced to fail. Moreover, no large cardinal axiom can
come to the rescue and decide the matter. The generic multiverse was
introduced as a means of taking such deep incompleteness seriously
as something we might just have to live with. In its simplest form,
a generic multiverse $\mathbb{V}_{M}$ is formed by taking a countable
transitive model $M$ of $ZFC$ and then considering all the models
that can be obtained from $M$ by finite sequences of generic extensions
and refinements \citep{WoodinGM}. We might then live with the incompleteness
of $ZFC$ by only taking seriously those statements that are true
in every world $N$ in $\mathbb{V}_{M}$. In particular, the continuum
hypothesis is relegated to the liminal space between always true and
always false. Woodin describes the motivation behind the generic multiverse
particularly clearly:
\begin{quote}
Is the generic-multiverse position a reasonable one? The refinements
of Cohen's method of \emph{forcing} in the decades since his initial
discovery of the method and the resulting plethora of problems shown
to be unsolvable, have in a practical sense almost compelled one to
adopt the generic-multiverse position.\footnote{I should mention at this point that although Woodin clearly describes
why embracing the generic multiverse might be attractive and understandable,
his goal in that paper is to argue against its adoption.} \citep{WoodinGM}
\end{quote}
If our goal is to take deep incompleteness seriously, then Woodin's
$\mathbb{V}_{M}$ has some limitations. $\mathbb{V}_{M}$ is a model
that we describe, so to speak, from the outside rather than foundational
framework like $ZFC$ than we can use, so to speak, from the inside.
It's still a valuable tool for understanding how the generic multiverse
position works, but it doesn't offer us a way of truly inhabiting
its multiverse. To address this, Steel introduced a computably axiomatizable
theory of the generic multiverse that is designed to be an autonomous
foundation for mathematics, which is on a par with $ZFC$, and within
which we can live with deep incompleteness \citep{SteelGP}.

Steel's theory has received a healthy dose of high quality philosophical
attention since its inception \citep{MeadMadPGMV,Meadows2Args,BAGARIA_TERNULLO_2023,BLUE_2024}.
My goals in this section and the rest of this paper are more logical
in spirit. I want to give an overview of: what $MV$ looks like from
the outside, through a rough classification of its models; and what
$MV$ feels like on the inside, by highlighting some significant consequences
of the theory. My thought is that by delivering a more comprehensive
overview of $MV$, the internal categoricity results below will make
more sense to the reader and perhaps over-optimistically that these
results might make $MV$ more accessible to a wider class of readers. 

\subsection{Notation}

Before we launch into the theory, let's first make some remarks about
notation. In general, I am aiming to follow the notation conventions
used by contemporary set theorists as used in, for example \citep{JechST,KunenST2}
and \citep{foreman2009handbook}. However, I'll highlight a few local
idiosyncrasies that I've taken up in order to make things look a little
smoother on the page. 

First and contrary to the usual conventions in mathematical logic,
a theory will just be a set of sentences that we usually call axioms.
We shall not demand, unless stated otherwise, that a theory is closed
under logical consequence. Given that internal categoricity is sensitive
to the axioms used to generate a theory, this is particularly important.\footnote{We should note that this approach to theories is quite common in the
theory of interpretability. See, for example, \citep{lindstrm2003aspects}
and \citep{EnayatLelykFOCat}.}

We shall, below, be frequently concerned with models of $\mathcal{L}_{\in}=\{\in\}$.
Most often these will be transitive models and for those, we adopt
the usual convention of denoting the model $\langle M,\in\rangle$
by its domain $M$. However, when we come to denote models that have
a non-standard membership relation we shall use the calligraphic font
and write $\mathcal{M}=\langle M,\in_{\mathcal{M}}\rangle$.

In almost all of the many simple forcing arguments that follow we
shall make great use of L\'{e}vy's collapse forcing. Recall that $Col(\omega,X)$
for $X\subseteq Ord$ is the partial order whose conditions are finite
partial functions $p:\omega\times X\rightharpoondown X$ where for
all $n\in\omega$ and $\alpha\in X$, $p(n,\alpha)<\alpha$ if it
is defined. We then order $Col(\omega,X)$ be reverse inclusion. Rather
than write $Col(\omega,X)$ again and again, we shall instead write
$\mathbb{C}_{X}$ where the $\mathbb{C}$ is intended to be mnemonic
for \uline{c}ollapse. We shall also use interval notation for sets
of ordinals so, for example $[\alpha,\beta)=\{\gamma\in\beta\ |\ \gamma\geq\alpha\}$.
If $H$ is $\mathbb{C}_{\beta}$-generic over $V$ for some $\beta\in Ord\cup\{Ord\}$
and $\alpha<\beta$, we write $H\restriction\alpha$ to denote $\{p\in\mathbb{C}_{\alpha}\ |\ p\in H\}$. 

The following theorem of Laver and Woodin plays a pivotal role in
the theory of the generic multiverse. We shall also use the terms
``generic refinement'' and ``ground'' as synonyms below. $U$
is a generic refinement of $W$ if there is some $U$-generic $G$
such that $W=U[G]$.
\begin{fact}
(Laver, Woodin)\label{fact:Laver,-Woodin} Every generic refinement
$U$ of $V$ is uniformly definable in $V$ using a parameter from
$U$. That is, there is a formula $\psi(x,y)$ of $\mathcal{L}_{\in}$
such that if $V=U[G]$ is a forcing extension of $U$ by a $U$-generic
filter $G\subseteq\mathbb{P}\in U$, then there is some $r\in U$
such that\footnote{See \citep{ReitzTGA} for a gentle proof and \citep{WoodinGM} for
something more general. } 
\[
U=\{x\ |\ \psi(x,r)\}.
\]
\end{fact}

Let us fix a formula $\psi(y,x)$ as in the statement of the fact
above. We shall write $x\in\mathsf{W}_{r}$ to mean that $\psi(x,r)$
holds. This is usually just written as $W_{r}$ rather than $\mathsf{W}_{r}$.
However, when talking about the generic multiverse and using the theory,
$MV$, upper case variables like $W$ are in high demand. Thus, we
use the sans serif font $\mathsf{W}$ to avoid confusion. 

\subsection{The $MV$ theory\label{subsec:The--theory}}

We are now ready to describe $MV$. The theory is articulated in the
language $\mathcal{L}_{MV}$ which is a two sorted language extending
$\mathcal{L}_{\in}$ with a new sort for \emph{worlds} beyond the
existing sort for \emph{sets}. Sets can be members of both sets and
worlds, but worlds can be members of neither sets nor worlds.\footnote{Strictly we should have two membership relations: one for when sets
are members of sets; and the other for when sets are members of worlds.
However, using just one symbol lightens our notation load in a helpful
way so we adopt this abuse of notation. } The axioms of $MV$ can then be stated somewhat formally as follows.
We shall write $gen(\mathbb{P},W)$ to denote the set of $\mathbb{P}$-generic
filters of a world $W$. 
\begin{lyxlist}{00.00.0000}
\item [{(World-extensionality)}] $\forall x(x\in W\leftrightarrow x\in U)\to W=U$
\item [{(World-transitivity)}] $x\in W\to x\subseteq W$
\item [{(World-ordinals)}] $\alpha\in Ord^{W}\leftrightarrow\alpha\in Ord^{U}$
\item [{($ZFC$-worlds)}] $\varphi^{W}$ where $\varphi$ is an axiom of
$ZFC$. 
\item [{(Set-capture)}] $\exists W\ x\in V$
\item [{(Extension)}] $\mathbb{P}\in W\exists U\exists g\in U\ (g\in gen(\mathbb{P},W)\wedge U=W[g])^{U}$
\item [{(Refinement)}] $r\in W\exists U\ (U=(\mathsf{W}_{r})^{W}$
\item [{(Amalgamation)}] 
\begin{align*}
 & \exists\mathbb{P}_{0}\in W_{0}\exists\mathbb{P}_{1}\in W_{1}\exists U\exists g_{0},g_{1}\in U\ \\
 & (g_{0}\in gen(\mathbb{P}_{0},W_{0})\wedge g_{1}\in gen(\mathbb{P}_{1},W_{1})\wedge\\
 & W_{0}[g_{0}]=U=W_{1}[g_{1}])^{U}.
\end{align*}
\end{lyxlist}
The first five axioms essentially say that worlds are inner models
of $ZFC$ and that every set $x$ is captured by a world $W$ in the
sense that $x\in W$. The final three axioms breathe life into the
\emph{generic} aspect of this multiverse. Extension tells us that
any generic extension of a world is itself a world. Refinement tells
us that any generic refinement of a world is a world. Amalgamation
tells us that any pair of worlds have a common generic extension. 

Extension and Refinement are straightforwardly motivated. They tell
us that the multiverse is closed under forcing and its inverse. Amalgamation
is a little more complicated.\footnote{See \citet{MeadMadPGMV} for some extensive discussion of this. }
One quite simple-minded motivation for including it is that it ensures
that our multiverse doesn't contain, so to speak, isolated galaxies
since any pair of worlds are subsumed into a larger world. However,
we shall see very soon that Amalgamation also allows us to obtain
particularly elegant models of $MV$ that have a structure which is
very familiar to set theorists. 

It's easy enough to offer these informal explanations but they only
get us so far. In the next section, we'll start thinking more seriously
about models of $MV$, but before we do this, let's ask a natural
question about the consequences of this theory: 
\begin{quote}
Given that $MV$ is a theory of sets and worlds, we can consider its
reduction to the set domain, which gives us a theory in $\mathcal{L}_{\in}$.
What is that theory like?
\end{quote}
It turns out that it has a very nice theory indeed. Informally, we
might say that it satisfies countable set theory. More precisely,
let $ZFC^{-}$ be $ZFC$ articulated using the Collection schema (rather
than Replacement) and a Set Induction schema (rather than Foundation)
where we remove the Powerset axiom.\footnote{Nothing important swings on using Set Induction rather than Foundation.
Like Collection, it just behaves better in weakenings of $ZFC$.} Let $ZFC_{count}^{-}$ be $ZFC^{-}$ with the following additional
axiom: $\forall x\ |x|\leq\omega$. This is the natural theory of
the hereditarily countable sets and it is implied by $MV$.\footnote{Note that Choice is already implied by $\forall x\ |x|\leq\omega$
but we retain the $C$ just to emphasize that Choice is available. }
\begin{thm}
$MV\vdash ZFC_{count}^{-}$.\label{thm:MV|-ZFC_count}
\end{thm}

Before we remark on the proof of this theorem, this is a good moment
to return our attention to the continuum hypothesis and consider how
it should be understood in the context of $MV$. We remarked above
that $MV$ was developed with an eye to taking the independence of
the continuum hypothesis seriously. But how is this achieved? It turns
out that there are two ways of thinking about this. On the first way
of understanding things, we treat each world in the multiverse as
a genuine alternative and we see then that there will be worlds where
the continuum hypothesis is true and worlds where it is false. This
is a theorem of $MV$ that is essentially a consequence of Cohen's
initial work on forcing \citep{Cohen1963-COHTIO-5}. We might call
this the \emph{local perspective}. From this point of view, the cardinality
of the continuum is not determined since it has many possible values
across the multiverse. By contrast, the \emph{global perspective}
takes up the point of view considered in Theorem \ref{thm:MV|-ZFC_count}.
Rather than worrying about whether the continuum hypothesis is true
in one world or false in another, we ask whether it is true in the
underlying, global, universe of sets that sits behind each of the
worlds in the multiverse. The theorem above then tells us that every
set is countable. Unless we wish to contradict Cantor's famous theorem,
this means that there cannot be a set of all subsets of the natural
numbers: $\mathcal{P}(\omega)$ does not exist. Thus, from the global
perspective the continuum hypothesis is simply inexpressible and thus,
arguably enjoys an even greater degree of indeterminacy than is exhibited
locally. I do not wish to put forward a claim that one perspective
is better than the other or that a choice must be made. After all,
they both fit the formal theory quite well. Regardless, it is clear
that -- either way -- the continuum hypothesis has no definitive
answer in $MV$. 

\subsection{The models of $MV$}

Let return to Theorem \ref{thm:MV|-ZFC_count} and its proof. It shouldn't
be too surprising to see that $MV$ implies every set is countable:
Extension ensures that for any cardinal, there will be a world where
it has been collapsed. Since Cantor's theorem tells us that any set
and its powerset have different cardinalities, we also see that the
Powerset Axiom must fail. To show that $ZFC^{-}$ is satisfied, it
will be helpful to consider some existing results about models of
$MV$. Let's begin with some technical preliminaries.

A model of $MV$ is of the form $\mathcal{W}=\langle\mathbb{W}^{\mathcal{W}},\mathbb{D}^{\mathcal{W}},\in^{\mathcal{W}}\rangle$
where $\mathbb{W}^{\mathcal{W}}$ is the world domain, $\mathbb{D}^{\mathcal{W}}$
is the world domain and $\in^{\mathcal{W}}$ is the membership relation.
In \citep{SteelGP} and \citep{MeadMadPGMV}, we are given a particularly
nice class of models for $MV$ that we might regard as its \emph{standard
models}. Let $\mathcal{M}=\langle M,\in^{\mathcal{M}}\rangle$ be
a model of $ZFC$ and suppose that $G$ is $(\mathbb{C}_{Ord})^{\mathcal{M}}$-generic
over $\mathcal{M}$. Then let $\mathcal{M}[G]\supseteq\mathcal{M}$
be the (class) generic extension of $\mathcal{M}$ via $G$.\footnote{If $\mathcal{M}$ isn't transitive, it's a little fiddly to ensure
that $\mathcal{M}\subseteq\mathcal{M}[G]$. A way of doing this is
described in Appendix A.1 of \citep{MeadMadPGMV}, where it is denoted
as $\mathcal{M}_{ult}[G]$, however, we'll gloss over that subtlety
here.} Note that since $(\mathbb{C}_{Ord})^{\mathcal{M}}$ is a pretame
forcing over $\mathcal{M}$, we see that $\mathcal{M}[G]$ satisfies
$ZFC^{-}$. A standard model $\mathcal{M}^{G}$ is then formed from
the set $\mathbb{W}^{\mathcal{M}^{G}}$ of generic refinements of
$\mathcal{M}[G\restriction\alpha]$ for each $\alpha\in Ord^{\mathcal{M}}$.\footnote{Or a little more precisely, $\mathcal{N}=\langle N,\in^{\mathcal{N}}\rangle$
is a world in $\mathcal{M}^{G}$ is there is some $\alpha\in Ord^{\mathcal{M}}$
and $r\in\mathcal{M}[G\restriction\alpha]$ such that $x\in N$ iff
$\langle\mathcal{M}[G],\mathcal{M},\in^{\mathcal{M}[G]}\rangle\models x\in(\mathbb{W}_{r})^{\mathcal{M}[G\restriction\alpha]}$. } Thus, $\mathcal{M}^{G}=\langle\mathbb{W}^{\mathcal{M}^{G}},\mathcal{M}[G],\in^{\mathcal{M}[G]}$.
We then recall that:
\begin{fact}
(Steel, Theorem 26 and 27 in \citep{MeadMadPGMV}) (1) If $\mathcal{M}$
is a model of $ZFC$ and $G$ is $(\mathbb{C}_{Ord})^{\mathcal{M}}$-generic
over $\mathcal{M}$, then $\mathcal{M}^{G}\models MV$.\label{fact:Steel,sound-comp}

(2) If $\mathcal{W}$ is a countable model of $MV$ and $\mathcal{M}$
is a world in $\mathcal{W}$, then for some $G$ that is $(\mathbb{C}_{Ord})^{\mathcal{M}}$-generic
over $\mathcal{M}$, we have\label{fact:Steel comp}
\[
\mathcal{W}=\mathcal{M}^{G}.
\]
\end{fact}

Informally speaking, given any model $\mathcal{M}$ of $ZFC$ and
a generic to collapse its ordinals, we may use them to obtain a standard
model of $MV$. And in the other direction, we see that if $\mathcal{W}$
is a \emph{countable }model of $MV$ then it must be a standard model.
We can use this to prove the theorem about the set part of $MV$. 
\begin{proof}
(of Theorem \ref{thm:MV|-ZFC_count}). Let $\mathcal{W}$ be a model
of $MV$. We'd like to apply Fact \ref{fact:Steel comp}(2), but for
that we need $\mathcal{W}$ to be countable. But if $\mathcal{W}$
is not countable, we can simply take a countable Skolem hull. Now
let $\mathcal{M}$ be an arbitrary world in $\mathcal{W}$. Fact \ref{fact:Steel comp}(2)
then tells us that we may fix some $(\mathbb{C}_{Ord})^{\mathcal{M}}$-generic
over $\mathcal{M}$ such that $\mathcal{W}=\mathcal{M}^{G}$. Then
since $(\mathbb{C}_{Ord})^{\mathcal{M}}$ is pretame, the set domain
$\mathcal{M}[G]$ of $\mathcal{M}^{G}$ satisfies $ZFC^{-}$. Moreover,
Extension ensures that every set is countable. 
\end{proof}
The Fact above is stated quite generally in that it accommodates
the worlds that are nonstandard in the sense that their membership
relation is not well-founded. But most (although not all) of the time,
we'll be dealing with models in which every world is transitive. In
those cases, we know that $\in^{\mathcal{W}}$ is just the \emph{real
}$\in$ restricted to $\mathbb{D}^{\mathcal{W}}$. Moreover, in these
case we can, without loss of generality, think of $\mathbb{W}^{\mathcal{W}}$
as literally being the set of those transitive sets. If we do this
then we see that $\mathbb{D}^{\mathcal{W}}=\bigcup\mathcal{W}$.\footnote{We'll see later that the converse does not hold: we cannot recover
the world domain from the set domain.} Thus, when we deal with transitive models of $MV$, we see that $\mathbb{D}^{\mathcal{W}}$
and $\in^{\mathcal{W}}$ are redundant and so in such cases, we'll
often abuse notation and let $\mathcal{W}$ just be its world domain
$\mathbb{W}^{\mathcal{W}}$. Now if we restrict our attention countable
transitive models we get a particularly nice model-theoretic representation
$MV$. 
\begin{cor}
For countable $\mathcal{L}_{MV}$-structures $\mathcal{W}$ containing
transitive worlds, the following are equivalent:\label{cor:Steel}
\begin{enumerate}
\item $\mathcal{W}\models MV$; and
\item $\mathcal{W}=M^{G}$ for some $G$ that is $Col(\omega,<Ord^{M})$-generic
over $M.$
\end{enumerate}
\end{cor}

For a helpful way to think about standard models of $MV$, it is helpful
to recall the proof behind Solovay's famous model where the reals
are Lebesgue measurable, have the property of Baire and the perfect
set property \citep{Solovay1970AMO}. We begin that proof by supposing
that there is an inaccessible cardinal $\kappa$ and then consider
a model where all the ordinals below $\kappa$ have been collapsed
using a generic $G$ for $\mathbb{C}_{\kappa}$. This was the first
of many similar models obtained by collapsing large cardinals. It
is a familiar playground for set theorists and -- as we'll see many
of the techniques for working there translate quite directly to models
of $MV$. The main difference is that to make a standard model of
MV we don't just collapse some of the ordinals, we collapse all of
them. As we've seen this gives us our set domain. We might then regard
the worlds of the multiverse are recording the history of all possible
forcings along the road to this great collapse.

Standard models arguably provide the most natural tool for understanding
$MV$ and how it is connected to $ZFC$. Moreover, this bridge provides
a helpful and well understood toolbox for proving theorems about $MV$.
However, not all models of $MV$ are standard. Given that we concerned
in this paper with questions around categoricity, it will be helpful
to survey the landscape beyond the realm of standard models. Of course,
$MV$ is just a two-sorted theory in first order logic and so, by
the upward L\"{o}wenheim-Skolem theorem, it is obvious that there are
models of $MV$ of every possible cardinality. But if we place some
natural constraints on the models and the worlds they contain a more
interesting picture emerges. For example, given Corollary \ref{cor:Steel}
we might wonder if every model $\mathcal{W}$ of $MV$ containing
countable transitive models must itself be countable. This is not
the case. 
\begin{prop}
If $MA(\kappa)$ holds for $\kappa\geq\omega$, then there are models
$\mathcal{W}$ of $MV$ with $|\mathcal{W}|=\kappa^{+}$ where every
world $M$ in $\mathcal{W}$ is countable and transitive (if there
is a countable transitive model of $ZFC$). Moreover, any such $\mathcal{W}$
will be nonstandard. 
\end{prop}

\begin{proof}
We start by observing that such models cannot be standard. To see
this, suppose $M$ is a world in $\mathcal{W}$ and $G$ is $\mathbb{C}_{Ord}^{M}$-generic
over $M$. Then for all $\alpha\in Ord^{M}<\omega_{1}$, $M[G\restriction\alpha]$
is countable and so it has countably many generic refinements, $\mathsf{W}_{r}^{M[G\restriction\alpha]}$,
indexed by elements $r\in M[G\restriction\alpha${]}. Thus, the worlds
of $M^{G}$ are a countable union of countable sets, while $\mathcal{W}$
is uncountable.

Now we prove that such models exist. Let $M$ be a countable transitive
model of $ZFC$. Before we build our target model, first observe that
if $G_{0}$ and $G_{1}$ are $\mathbb{C}_{Ord^{M}}$-generic over
$M$, then we can easily find $G_{2}$ that is $\mathbb{C}_{Ord^{M}}$-generic
over $M$ where $M[G_{2}]=M[G_{0}\times G_{1}]$. Thus, we obtain
a standard model $M^{G_{2}}$ extending both $M^{G_{0}}$ and $M^{G_{1}}$.
Clearly, this can be generalized to any finite set of $\mathbb{C}_{Ord^{M}}$-generics.
We might say that any finite collection standard models based on $M$
can be amalgamated into another standard model. Similar remarks apply
to countable collections of standard models, but as we've seen no
amalgamation fact can hold for uncountable collections of standard
models. We address this problem by inductively defining a sequence
$\langle G_{\alpha}\rangle_{\alpha<\kappa^{+}}$ of filters that are
mutually $\mathbb{C}_{Ord^{M}}$-generic over $M$. Let $\alpha<\kappa^{+}$
and suppose that we've already defined $\langle G_{\beta}\rangle_{\beta<\alpha}$
such that for all finite $X\subseteq\alpha$, $\prod_{\gamma\in X}G_{\gamma}$
is $\prod_{\gamma\in X}\mathbb{C}_{Ord^{M}}$-generic over $M$.\footnote{Note that $\prod_{\gamma\in X}\mathbb{C}_{Ord^{M}}$ is just $(\mathbb{C}_{Ord^{M}})^{|X|}$.}
We wish to define $G_{\alpha}$ so that $\langle G_{\beta}\rangle_{\beta<\alpha+1}$
has the analogous property at $\alpha+1$. For finite $X\subseteq\alpha$,
let 
\[
\mathcal{D}_{X}=\{D\in M[\prod_{\gamma\in X}G_{\gamma}]\ |\ D\text{ is dense in }\mathbb{C}_{Ord^{M}}\}.
\]
Then let $\mathcal{D}=\bigcup_{X\in[\alpha]^{<\omega}}\mathcal{D}_{X}$.
Since $\alpha<\kappa^{+}$, we see $|\mathcal{D}|\leq\kappa$ and
since $\mathbb{C}_{Ord^{M}}$ has the ccc (in $V$), we may use $MA(\kappa)$
we may fix a filter $G_{\alpha}$ where $G_{\alpha}\cap D\neq\emptyset$
for all $D\in\mathcal{D}$. This ensures that for all finite $X\subseteq\alpha$,
$G_{\alpha}$ is $\mathbb{C}_{Ord^{M}}$-generic over $M[\prod_{\gamma\in X}G_{\gamma}]$
and application of the product lemma completes our inductive argument.
Finally, we then let $\mathcal{W}$ be the set of generic refinements
of models of the form $M[\prod_{\gamma\in X}(G_{\gamma}\restriction\alpha)]$
where $X\in[\kappa^{+}]^{<\omega}$ and $\alpha\in Ord^{M}$. 
\end{proof}
Thus, we see that in contexts where some form of Martin's axiom holds,
there will be large nonstandard models of $MV$ in which every world
is countable and transitive. Moreover, since $MA(\omega)$ is a theorem
of $ZFC$, we see that such nonstandard models are unavoidable.
\begin{cor}
There are models $\mathcal{W}$ of $MV$ with $|\mathcal{W}|=\aleph_{1}$
where every world $\mathcal{M}$ in $\mathcal{W}$ is countable and
transitive. 
\end{cor}

Given that the delivery of generic sets for models of $ZFC$ is often
dependent on those models being countable, we might also wonder whether
standard models of $MV$ can only contain countable models. This is
not the case. 
\begin{prop}
If $0^{\#}$ exists, there there are standard models $\mathcal{W}$
of $MV$ that contain uncountable transitive worlds.\label{prop:Hjorth}
\end{prop}

\begin{proof}
Since $0^{\#}$ exists, we know that $L_{\omega_{1}}$ satisfies $ZFC$.
To deliver a standard model of $MV$ as desired, it will suffice to
show that there is some $G$ that is $\mathbb{C}_{\omega_{1}}$-generic
over $L_{\omega_{1}}$. Since $\mathbb{C}_{\omega_{1}}$ is not countable,
the usual Baire category argument will not work. We address this by
defining the required generic inductively.\footnote{A version of this argument can be found in the first claim within
the proof of Theorem 3.5 in \citep{HjorthU2}. I learned it from a
MathOverflow response by Farmer Schlutzenberg.}

First, we observe that for all $\alpha<\omega_{1}$, $\mathcal{P}(\mathbb{C}_{\alpha})^{L}$
is countable and so we may obtain a $\mathbb{C}_{\alpha}$-generic
filter over $L_{\omega_{1}}$ for any such $\alpha$. To set up the
induction, fix the club sequence $\langle\kappa_{\alpha}\rangle_{\alpha\leq\omega_{1}}$
of indiscernibles below $\omega_{1}$ witnessing that $0^{\#}$ exists
and where $\kappa_{\omega_{1}}=\omega_{1}$. We aim to deliver a sequence
$\langle G_{\alpha}\rangle_{\alpha\leq\omega_{1}}$ such that:
\begin{enumerate}
\item $G_{\alpha}$ is $\mathbb{C}_{\kappa_{\alpha}}$-generic over $L$;
and 
\item $G_{\alpha}\subseteq G_{\beta}$ for all $\alpha\leq\beta\leq\omega_{1}$. 
\end{enumerate}
For successor ordinals, $\alpha+1$, we just find $H$ that is $\mathbb{C}_{[\kappa_{\alpha},\kappa_{\alpha+1})}$-generic
over $L[G_{\alpha}]$ and then $G_{\alpha+1}$ can be what is essentially
the product of $G_{\alpha}$ and $H$. At limit ordinals $\beta$,
we have to be a little more careful. Note first that each $\kappa_{\beta}$
is inaccessible in $L$. Thus, $L$ thinks that $\mathbb{C}_{\kappa_{\beta}}$
has the $\kappa_{\beta}$-cc property and so any maximal antichain
$A$ in $\mathbb{C}_{\kappa_{\beta}}$ will be a subset of $\mathbb{C}_{\kappa_{\alpha}}$
for some $\kappa_{\alpha}<\kappa_{\beta}$.\footnote{See the first lemma on page 15 of \citep{Solovay1970AMO} for the
canonical proof of this fact.} Given this, we can let $G_{\beta}=\bigcup_{\alpha<\beta}G_{\alpha}$. 
\end{proof}
This result depends on a strong assumption beyond $ZFC$. Without
it, it is possible to find models $\mathcal{W}$ of $MV$ where every
world is transitive, but $\mathcal{W}$ is not standard.

\begin{prop}
If $L_{\omega_{1}}\models ZFC$, then it is consistent that there
is a model $\mathcal{W}$ of $MV$ where every world $M$ in $\mathcal{W}$
is uncountable and transitive, but $\mathcal{W}$ is not standard.
\end{prop}

\begin{proof}
Again, we are going to build a multiverse from $L_{\omega_{1}}$.
However, we know from Proposition \ref{prop:Hjorth} that this cannot
be done within a model where $0^{\#}$ exists. More precisely, we
need to find a model where there is a $\mathbb{C}_{\alpha}$-generic
set for every $\alpha<\omega_{1}$, but no $\mathbb{C}_{\omega_{1}}$-generic
set exists. To obtain such a model, we use Silver's collapse forcing.
For $X\subseteq Ord$, we let $\mathbb{S}_{X}$ be the poset whose
domain is formed from functions $p:s\times Y\to Ord$ where we have:
the \emph{first domain}, $s\in[\omega]^{<\omega}$; the \emph{second
domain},\emph{ }$Y\in[X]^{\leq\omega}$; and for all $n\in s$ and
$\alpha\in Y$, $p(n,\alpha)<\alpha$.\footnote{See Section 20 in Cummings' chapter in \citep{foreman2009handbook}
for more information on this and a more general definition.} Silver's collapse forcing can be obtained from L\'{e}vy's collapse by
weakening the constraint on the second domain so that it merely has
to be countable rather than finite. This gives $\mathbb{S}_{\omega_{1}}$
quite different properties from $\mathbb{C}_{\omega_{1}}$. In particular,
$\mathbb{S}_{\omega_{1}}$ has, so to speak, a certain amount of countable
completeness in its second domain. 

Now let $H$ be $\mathbb{S}_{\omega_{1}}$-generic over $L$.\footnote{Strictly speaking, there is no reason to think such an $H$ exists.
However following usual conventions, we can think of ourselves as
now working \emph{virtually} within the context of $\Vdash_{\mathbb{S}_{\omega_{1}}}$
as calculated in $L$. This is sufficient for our consistency claim.
Ordinarily, such metamathematical subtleties can be largely ignored.
But given that this paper is so concerned with generic sets and whether
they exist, it is worth taking the time to take a little care around
this.} Since $\mathcal{P}(\mathbb{C}_{\alpha})^{L}$ is countable for all
$\alpha<\omega_{1}$, we see that there is a $\mathbb{C}_{\alpha}$
-generic over $L$ for each such $\alpha$. It will thus suffice to
show that $L[H]$ contains no sets $G$ that are $\mathbb{C}_{\omega_{1}}$-generic
over $L$.\footnote{Our proof is an elaboration of a sketch by Mohammad Golshani on MathOverflow.
Nothing in the argument relies on us using $L$ rather than $V$.
It just lines up more smoothly with the story of the proof. I'm grateful
to Jason Chen for drawing my attention to it.} To demonstrate this, we prove a claim that suffices for the proposition.
\begin{claim*}
(1) If $A\in[\omega_{1}]^{\omega_{1}}\cap V[H]$, then there is a
countably infinite $X\subseteq A$ in $L$.

(2) If $G$ is $\mathbb{C}_{\omega_{1}}$-generic over $L$, then
for some $A\in[\omega_{1}]^{\omega_{1}}\cap L[G]$ there is no countably
infinite $X\subseteq A$ in $L$. 
\end{claim*}
To see that this is sufficient, suppose toward a contradiction that
$G\in V[H]$ is $\mathbb{C}_{\omega_{1}}$-generic over $L$. Then
by (1) we may fix $A\in[\kappa]^{\kappa}\cap V[G]$ which has no countably
infinite subset in $L$. But then we see that $A\in V[G]\subseteq V[H]$
and so by (2) $A$ must have a countable infinite subset, which is
impossible. 

To prove (1), we make an argument that exploits the version of countable
completeness possessed by the second domain of $\mathbb{S}_{\omega_{1}}$.\footnote{The argument is a simple generalization of the standard argument that
forcings closed under $\omega$-sequences don't add any new $\omega$-sequences.
See, for example, Theorem VII.6.14 in \citep{KunenST}.} We suppose toward a contradiction that $L$ contains no injection
from $\omega$ into $A$. We continue by letting $\dot{A}\in L^{\mathbb{S}_{\omega_{1}}}$
be a name such that $\dot{A}_{H}=A$ and fixing $p\in H$ that forces
the statement: there is is no injection $f:\omega\to\dot{A}$ with
$f\in\check{V}$. Working in $L$, we then define sequences $\langle p_{n}\rangle_{n\in\omega}$
and $\langle\alpha_{n}\rangle_{n\in\omega}$ such that $p_{0}\leq p$
and for all $n\in\omega$:
\begin{itemize}
\item $p_{n}\Vdash\alpha_{n}\in\dot{A}$;
\item $\alpha_{n+1}>\alpha_{n}$; 
\item $p_{n+1}\leq p_{n}$; and 
\item $dom(p_{n})=m\times Y$ where $m\in\omega$ is least among the second
domains of conditions satisfying the first three bullet points.
\end{itemize}
We claim that there is some $k\in\omega$ such that for all $n\in\omega$,
if $m$ is the second domain of $p_{n}$, then $m\leq k$. To see
this, we move back to $V$ and observe that since $|A|=\kappa$ and
there are only countably many possibilities for the first domains
of conditions, there must be some $k\in\omega$ such that there are
infinitely many $\alpha\in A$ with some $p\in H$ with $p\Vdash\alpha\in\dot{A}$
where the second domain of $p$ is $<\kappa$. Since we have a bound
on the first domain, it can be seen that $q=\bigcup_{n\in\omega}p_{n}\in\mathbb{S}_{\omega_{1}}$.
Finally, let $J$ be $\mathbb{S}_{\omega_{1}}$-generic over $L$
with $q\in J$. Then we see that $\langle\alpha_{n}\rangle_{n\in\omega}$
is an injection from $\omega$ into $A$ that is an element of $L$,
which gives a contradiction completing our proof of (1).

To prove (2), we provide a $\mathbb{C}_{\omega_{1}}$-name $\dot{A}$
for a set $A$ whose realization has no countably infinite subset
in $L$. More specifically, we let 
\[
\dot{A}=\{\langle\check{\alpha},p\rangle\in\check{\omega}_{1}\times\mathbb{C}_{\omega_{1}}\ |\ p(0,\alpha)=1\}
\]
and $A=\dot{A}_{G}$. The see the idea behind this name, let $g=\bigcup G:\omega\times\omega_{1}\to\omega_{1}$
and note that $\alpha\in\dot{A}_{G}$ iff $g(0,\alpha)=1$. So we
are restricting our attention to how conditions $p\in\mathbb{C}_{\omega_{1}}$
act on $0$ in the first domain and we only care about whether the
output is $1$ or not. This gives us a name for an element of $[\omega_{1}]^{\omega_{1}}\cap L[G]$.

Next we claim that $A$ is $Add(\omega_{1},1)$-generic over $V$.
To see this, we recall a basic, abstract forcing fact. Let us say
that $\sigma:\mathbb{P}\to\mathbb{Q}$ is \emph{weakly dense} if $\sigma$
is order preserving and for all $p\in\mathbb{P}$ and $q\in\mathbb{Q}$
with $q\leq\sigma(p)$, there is some $p^{*}\in\mathbb{P}$ such that
$p\not\perp p^{*}$ and $\sigma^{*}(p)\leq q$. It is well-known then
that if $J$ is $\mathbb{P}$-generic over $V$, then the upward closure
of $\sigma``J$ is $\mathbb{Q}$-generic over $V$.\footnote{See Lemma 15.45 in \citep{JechST} for a proof.}
Now let us construe $Add(\omega_{1},1)$ as the poset whose conditions
$q$ are finite partial functions from $\omega_{1}$ to $2$; and
let $\sigma:\mathbb{C}_{\omega_{1}}\to Add(\omega_{1},1)$ be such
that for all $p\in\mathbb{C}_{\omega_{1}}$, $\sigma(p)(\alpha)=1$
iff $p(0,\alpha)=1$.\footnote{This is not strictly correct. Since $p(0,0)$ has no possible output
and $p(0,1)$ can only output $0$, we can't just use $\alpha$ on
both sides. Fixing this makes a simple idea ugly on the page, so I'll
stick with the sloppiness. It is, however, easily addressed and we
leave that to the reader.} It can be seen that $\sigma$ is a weakly dense embedding and so
$J=\sigma``G$ is $Add(\omega_{1},1)$-generic over $L$. Moreover,
it is not difficult to see that $A=\bigcup J$. Finally, suppose toward
a contradiction that $X\in L$ and $X\subseteq A$. Now let $\dot{J}$
be the canonical name for a $Add(\omega_{1},1)$-generic set.\footnote{Here we mean that $\dot{J}=\{\langle\check{q},q\rangle\ |\ q\in\mathbb{Q}\}$.}
Then we may fix $q\in J$ such that $q\Vdash X\subseteq\bigcup\dot{J}$.
But this means that we have a finite condition $q\in Add(\omega_{1},1)$
determining infinitely many membership facts about $\bigcup\dot{J}$,
which is impossible.
\end{proof}

\begin{thebibliography}{45}
\providecommand{\natexlab}[1]{#1}
\providecommand{\url}[1]{\texttt{#1}}
\expandafter\ifx\csname urlstyle\endcsname\relax
  \providecommand{\doi}[1]{doi: #1}\else
  \providecommand{\doi}{doi: \begingroup \urlstyle{rm}\Url}\fi
\bibitem[Woodin(2012)]{WoodinGM} W.~Hugh Woodin.
\newblock \emph{The Continuum Hypothesis, the Generic Multiverse of Sets, and
  the $\Omega$ Conjecture}.
\newblock Cambridge University Press, 2012.
\bibitem[Steel(2014)]{SteelGP} John~R. Steel.
\newblock G{\"o}del's program.
\newblock In Juliette Kennedy, editor, \emph{Interpreting G\"o{}del: Critical
  Essays}. Cambridge University Press, 2014.
\bibitem[Maddy and Meadows(2020)]{MeadMadPGMV} Penelope Maddy and Toby Meadows.
\newblock A reconstruction of {S}teel's multiverse project.
\newblock \emph{Bulletin of Symbolic Logic}, 26\penalty0 (2):\penalty0
  118--169, 2020.
\bibitem[Meadows(2021)]{Meadows2Args} Toby Meadows.
\newblock Two arguments against the generic multiverse.
\newblock \emph{Review of Symbolic Logic}, 14\penalty0 (2):\penalty0 347--379,
  2021.
\bibitem[Bagaria and Ternullo(2023)]{BAGARIA_TERNULLO_2023} Joan Bagaria and Claudio Ternullo.
\newblock Steel's programme: Evidential framework, the core and ultimate-{L}.
\newblock \emph{The Review of Symbolic Logic}, 16\penalty0 (3):\penalty0
  788--812, 2023.
\bibitem[Blue(2024)]{BLUE_2024} Douglas Blue.
\newblock The generic multiverse is not going away.
\newblock \emph{The Review of Symbolic Logic}, pages 1--33, 2024.
\newblock \doi{10.1017/S1755020324000297}.
\bibitem[Jech(2003)]{JechST} Thomas Jech.
\newblock \emph{Set Theory}.
\newblock Springer, Heidelberg, 2003.
\bibitem[Kunen(2011)]{KunenST2} Kenneth Kunen.
\newblock \emph{Set Theory}.
\newblock College Publications, London, 2nd edition, 2011.
\bibitem[Foreman and Kanamori(2009)]{foreman2009handbook} M.~Foreman and A.~Kanamori.
\newblock \emph{Handbook of Set Theory}.
\newblock Springer Netherlands, 2009.
\bibitem[Lindstr{\"o}m(2003)]{lindstrm2003aspects} P.~Lindstr{\"o}m.
\newblock \emph{Aspects of Incompleteness: Lecture Notes in Logic 10}.
\newblock Lecture notes in logic. Taylor \& Francis, 2003.
\bibitem[Enayat and \L{}e\l{}yk(2024)]{EnayatLelykFOCat} Ali Enayat and Mateusz \L{}e\l{}yk.
\newblock Categoricity-like properties in the first order realm.
\newblock \emph{Journal for the Philosophy of Mathematics}, 1:\penalty0 63--98,
  Sep. 2024.
\bibitem[Reitz(2007)]{ReitzTGA} Jonas Reitz.
\newblock The ground axiom.
\newblock \emph{Journal of Symbolic Logic}, 72\penalty0 (4):\penalty0
  1299--1317, 2007.
\bibitem[Cohen(1963)]{Cohen1963-COHTIO-5} Paul Cohen.
\newblock The independence of the continuum hypothesis.
\newblock \emph{Proc. Nat. Acad. Sci. USA}, 50\penalty0 (6):\penalty0
  1143--1148, 1963.
\bibitem[Solovay(1970)]{Solovay1970AMO} Robert~M. Solovay.
\newblock A model of set-theory in which every set of reals is {L}ebesgue
  measurable*.
\newblock \emph{Annals of Mathematics}, 92:\penalty0 1--56, 1970.
\bibitem[Hjorth(1995)]{HjorthU2} Greg Hjorth.
\newblock The size of the ordinal u2.
\newblock \emph{Journal of the London Mathematical Society}, 52\penalty0
  (3):\penalty0 417--433, 1995.
\bibitem[Kunen(2006)]{KunenST} Kenneth Kunen.
\newblock \emph{Set Theory: an introduction to independence proofs}.
\newblock Elsevier, Sydney, 2006.
\end{thebibliography}
\end{document}
